\chapter{Первые двенадцать учебных недель}

\section{Две недели диагностики и восстановления}
В первые две недели выбрать четыре из восьми задач ниже, обязательно взяв
одну топологическую и одну аналитическую. На каждую выделить до 90 минут
первой самостоятельной попытки, затем до двух часов разбора и исправления.
Остальные задачи остаются ориентирами на квартал. Весь список не требуется
решить за две недели.

\begin{exercise}[Локальная связность и гипотезы]
Дать определение локальной связности подпространства и критерий через
связные окрестности внутри произвольной окрестности. Привести два связных
компактных множества в плоскости с несвязным непустым пересечением.
Объяснить, почему принадлежность центральной точки пересечению не помогает.
\end{exercise}

\begin{exercise}[Критическая орбита]
Доказать, что ограниченность критической орбиты \(f_c(z)=z^2+c\)
влечёт \(|c|\le2\) и \(|f_c^n(0)|\le2\) для всех \(n\).
Отсюда вывести компактность \(\M\) через непрерывность конечных орбитальных
полиномов. Если начинается геометрический рост, указать его постоянный
множитель, больший единицы.
\end{exercise}

\begin{exercise}[Настоящий модуль]
Для \(0<r<R\) вычислить экстремальную длину семейства кривых,
соединяющих две граничные окружности кольца \(A(r,R)\).
Получить \(\modu A=(2\pi)^{-1}\log(R/r)\).
Сначала объявить нормировку \(\sup_\rho L_\rho(\Gamma)^2/A_\rho\),
затем доказать верхнюю и нижнюю оценки. Проверить масштабирование.
\end{exercise}

\begin{exercise}[Компактные вложенные множества]
Пусть непустые компакты \(K_n\) в метрическом пространстве убывают и
\(\bigcap K_n=\{x\}\). Доказать \(\diam K_n\to0\).
Отдельно разобрать, где используется компактность \(K_0\), и построить
контрпример без неё. Версия с окрестностями: для любой открытой \(U\ni x\)
найдётся \(n\), для которого \(K_n\subset U\).
\end{exercise}

\begin{exercise}[Лемма Шварца--Пика]
Нормировать автоморфизмами диска произвольное голоморфное
\(h:\D\to\D\) и вывести
\[
\frac{|h'(z)|}{1-|h(z)|^2}\le\frac1{1-|z|^2}.
\]
Разобрать случай равенства. Указать, что неувеличение гиперболического
расстояния ещё не даёт равномерного сжатия с постоянной меньше единицы.
\end{exercise}

\begin{exercise}[От колец к сжатию]
Сформулировать последовательное неравенство Грётша с точными условиями
взаимного расположения колец. Вывести один критерий сжатия вложенных
топологических дисков, явно указав нужную теорему о конформном модуле.
Проверить аргумент на концентрических дисках и сравнить с гауссовой площадью.
На диагностике достаточно локализовать недостающую аналитическую лемму;
её полное освоение запланировано отдельно.
\end{exercise}

\begin{exercise}[Кванторы условного корня]
Развернуть \repo{UniformOuterBuffer} на бумаге и воспроизвести из него
связную окрестность точки предельного компакта. Не использовать корневую
теорему как посылку. Отметить, что изменится, если \(\delta\) зависит от \(N\).
\end{exercise}

\begin{exercise}[Алгебраическая ветка]
Дать определения нижнего центрального ряда \(\gamma_n(G)\) и размерного
ряда относительно явно выбранного коэффициентного кольца. Для свободной
абелевой группы вычислить оба ряда. Изобразить, как подъём петли при
\(z\mapsto z^2:\C^*\to\C^*\) действует на фундаментальной группе.
Объяснить, какие дополнительные данные нужны для перехода к бисету.
\end{exercise}

\criterion{Если выбранные базовые задачи восстанавливаются самостоятельно,
не повторять весь вводный курс: переходить к новым разделам и использовать
освободившиеся сессии для исследовательских примеров. Если возникли пробелы,
вставить двухнедельный ремонт конкретной зависимости, не начинать книгу сначала.}

\section{Расписание оставшихся недель}
Подробные действия и файлы результатов перечислены в
\repo{first_12_weeks.md}. Здесь указана математическая логика.

\begin{table}[!tp]
\centering\small
\caption{Первые двенадцать учебных недель}
\begin{tabularx}{\textwidth}{@{}P{0.07\textwidth}P{0.37\textwidth}Y@{}}
\toprule
Нед. & Материал & Один основной результат\\
\midrule
1 & Диагностика: орбиты и топология & Две попытки без подсказок; одна
исправленная карточка с доказательством\\
2 & Аудит определений и две оставшиеся диагностические задачи & Сопоставление
модуля кольца и weighted area; карта условного корня\\
3 & Любич, §6, с.~119--132: выбранные пункты & Полное вычисление модуля
круглого кольца и проверка нормировок\\
4 & Последовательное неравенство и вложенные компакты & Одно доказательство
сжатия в явно указанной общности; квартальный журнал ещё не нужен, только
первый четырёхнедельный отчёт\\
5 & Любич, §7, с.~133--145: начало & Вывод метрики диска выбранной
нормировки и инвариантности относительно автоморфизмов\\
6 & Продолжение §7 & Доказательство Шварца--Пика, случай равенства,
пример отсутствия равномерной строгой константы\\
7 & Компактность и компоненты из текущего проекта & Восстановление одного
перехода в доказательстве достаточности буфера\\
8 & Любич, §8, с.~146--151 & Точная формулировка граничного продолжения
Риманова отображения; пример, где продолжение не инъективно\\
9 & Любич, приложение §10, с.~163--172; выбранные леммы & Принцип минимума
для гармонической функции в нужной ограниченной области; условия на границе\\
10 & Вводный Бёттхер: §23, с.~312--329, один фрагмент & Функциональное
уравнение Green-функции из определения и нулевой уровень на ограниченной орбите\\
11 & Любич, §9, с.~152--162; пример замороженной модели & Сравнение полного
Green-подуровня и граф-разрезанного puzzle; конечный пример \(c=0\) или \(-1\)\\
12 & Повторение и письменная сборка & Связная записка с 3--5 проверенными
утверждениями, устный доклад и пересмотр паспорта\\
\bottomrule
\end{tabularx}
\end{table}

Страницы относятся к местной книге 2023 года. Порядок намеренно не
совпадает с последовательным чтением всех страниц: пропускаемые разделы
возвращаются тогда, когда понадобятся их результаты. В неделе~10 не
требуется за две сессии доказать весь Бёттхер и глобальную униформизацию.

\section{Первые законченные тексты}
К концу квартала ожидаются примерно 6--8 небольших проверенных доказательств,
3--5 разобранных контрпримеров и одна объединяющая записка ориентировочно на
8--15 страниц. Это ориентиры для масштаба, а не норматив объёма. Если одна
лемма потребовала три недели и полностью понята, она ценнее десятка
непроверенных «готовых» доказательств.

В записке удобно связать три темы: настоящий модуль кольца; компактное
сжатие и компоненты; граница перехода от динамики к параметрам. Крупные
внешние теоремы перечисляются явно. Записка проверяет владение материалом,
а не заявляет новый результат.

\section{Первое внешнее обсуждение}
В недели~4--8 подготовить одну страницу о своей подготовке и двух возможных
темах, приложить 2--4 страницы одной леммы. Самостоятельно обратиться к
исследователю или семинару с конкретным вопросом о посильной задаче.
Цель --- получить содержательную критику проверяемого аргумента.
